\subsubsection{文件的打开与关闭}
在src文件夹中的fd.c文件中实现对文件的打开与关闭,这里打开与关闭文件是便于后续对文件的读写操作,具体代码如下:
\begin{ccode}{文件的打开与关闭}
#include "fat.h"

int fs_open(const char *filename,int mode){
    if (filename == NULL) return -1;

    //在当前目录下找普通文件
    int  dir_idx = -1;
    for(int i = 0; i < MAX_ENTRY; i++){
        if(Directory[i].ac == 1 &&
            Directory[i].parent == current_dir &&
            Directory[i].type == 0&&
            strcmp(Directory[i].name,filename) == 0
        ){
            dir_idx = i;
            break;
        }
    }
    if (dir_idx == -1){
        //没有找到
        printf("open failed : file not found : %s \n",filename);
        return -1;
    }

    // 找空闲fd
    for (int fd = 0; fd < MAX_FD; fd ++){
        if(FDTable[fd].used == 0){
            FDTable[fd].used = 1;
            FDTable[fd].dir_index = dir_idx;
            FDTable[fd].offset = 0;
            FDTable[fd].mode = mode;
            printf("open success : %s,fd = %d\n",filename,fd);
            return fd;
        }  
    }
    printf("open failed:too many open files\n");
    return -1;
}

bool fs_close(int fd){
    //无效fd
    if (fd < 0 || fd > MAX_FD){
        printf("close failed : invalid fd\n");
        return false;
    }
    //文件还未打开
    if(FDTable[fd].used == 0){
        printf("close failed : fd is not open yet\n");
        return false;
    }
    //正常关闭文件
    FDTable[fd].used = 0;
    FDTable[fd].dir_index = -1;
    FDTable[fd].offset = 0;
    FDTable[fd].mode = 0;
    printf("close success: fd = %d\n", fd);
    return true;
}

bool fs_seek(int fd,int offset){
    if(fd < 0 || fd >= MAX_FD) return false;
    if(FDTable[fd].used == 0) return false;
    if(offset < 0) return false;

    int dir_idx = FDTable[fd].dir_index;
    int file_size = Directory[dir_idx].size;

    // 读模式不允许超过文件大小
    if(FDTable[fd].mode == MODE_READ || FDTable[fd].mode == MODE_RDWR){
        if(offset > file_size){
            printf("seek failed: offset %d exceeds file size %d\n", offset, file_size);
            return false;
        }
    }

    FDTable[fd].offset = offset;
    return true;
}
\end{ccode}
其中fs\_open()函数用于打开当前目录下指定文件名的文件(这里为了简化实现,没有实现打开任意目录下的文件),通过遍历Directory数组,找到不为空、文件名为指定文件且父目录为当前目录的文件即为要打开的文件,找不到就返回失败并打印错误信息。
找到后再在FDTable数组(即打开文件表)中查找未使用的项,将其信息同步为当前要打开的文件,并将FDTable数组中对应项的索引返回(即文件描述符),后续对文件的读写操作都基于该索引。fs\_close()函数用于关闭指定文件,因为确定只有确定已经打开的文件
才需要关闭,因此这里关闭文件函数接收的参数为文件描述符fd。具体实现为检查fd未越界且FDTable数组中对应表项不为空时,将这一项清零即可。
fs\_seek()函数用于设置文件从何处开始读写,通过接收参数文件描述符fd和偏移量offset,在确认都没有越界的情况下更新FDTable中对应项的offset属性即可。
